\section{Uniformity for a small power of the conductor}
\label{sec:uniform-deformation}
The conductor can vary with the height after the errors in the moment
formulas and the deleted Euler factors are controlled uniformly.
The estimates below provide constants independent of the character
and conductor throughout the stated range.

\begin{theorem}\label{thm:uniform-deformation}
There exist absolute $\kappa,c,c_1,T_0>0$ such that, for $T\ge T_0$
and every primitive character of conductor $q\le T^\kappa$, each
strict initial normal-velocity sign has at least
$10^{-30}N_{\chi,T}$ distinct simple centers in $(T,2T)$.
Their same branches are simple, disjoint and have ordinates in
$(T,2T)$ throughout
\[
                  (1+T|u|)e^{12T|\Im u|}\le c_1e^{s_T},
              \qquad s_T=\sqrt{(\log\log T)/2}.
\]
The branches are holomorphic in the interior and continue locally
across its boundary. This region contains the complex disk
$|u|\le\sqrt{\log\log T}/(100T)$ and the real interval
$|u|\le c_1e^{s_T}/(2T)$ eventually. Their strict normal sides persist
for every positive real $u\le c/T$.

For each fixed finite interval $I$ of positive length, the lower limit
of the infimum over these characters of each signed counting
fraction in Theorem~\ref{thm:gaussian-speed-component} is at least
$10^{-30}\mathbb P(G\in I)$. For every fixed $A>0$ each initial
strict sign also has a possibly different set of at least
\[
               10^{-31}\mathbb P\{G\in[A+1,A+2]\}N_{\chi,T}
\]
centers with the same continuation and ordinate properties on
$(1+T|u|)e^{12T|\Im u|}\le c_Ae^{2As_T}$, and strict sides on
$0<u\le c'_A/T$. The constants and eventual thresholds in these
last assertions may depend on the fixed interval or on $A$, and are
uniform in the displayed conductor range.
\end{theorem}

\subsection{Uniform zero counts and moment errors}
\begin{proof}[Proof of Theorem~\ref{thm:uniform-deformation}]
Retain the single fixed $P,\nu$ used throughout and write
\[
 L=\log T,\quad V=\log L,\quad d=\tfrac12-2\nu>0,
 \qquad \ell=\frac{\log(qT)}L.
\]
All preliminary bounds below hold for $2\le q\le T$ and
$y\le T^\theta$, with fixed $\theta<1/2$.
The uniform Hadamard identity used to prove
\eqref{eq:uniform-character-growth}, evaluated at $2+it$, gives
\[
 \sum_\rho\frac{3/2}{(3/2)^2+(t-\gamma)^2}
            \ll\log(q(|t|+3)).
\]
Every term is positive. The local zero count is therefore
$O(\log(q(|t|+3)))$. Integration over $[T,2T]$ gives
\begin{equation}\label{eq:uniform-count-new}
 N_{\chi,T}=\frac T{2\pi}\{\log(qT)+O(1)\}
                              +O(\log(qT)\log T).
\end{equation}
Indeed the integrated gamma and conductor terms are
$\tfrac12T\log(qT)+O(T)$, and the Euler logarithmic derivative has
bounded integral by termwise integration of its absolutely convergent
series. The Poisson kernel integral differs from
$\pi\mathbf1_{(T,2T)}(\gamma)$ by at most
$C\{(1+|\gamma-T|)^{-1}+(1+|\gamma-2T|)^{-1}\}$ for
$|\gamma|\le4T$. The local count bounds this total error by
$O(\log(qT)\log T)$; dyadic tails give $O(\log(qT))$.
The same argument includes conductor one after adding the two
polynomial factors of its entire completion.

Each unit interval used for a top or bottom contour contains a height
at distance $\gg1/\log(q(|t|+3))$ from every ordinate. Subtracting
the completed partial fractions at $2+it$ from those at $\sigma+it$
gives $L_\chi'/L_\chi=O(\log^2(q(|t|+3)))$ there on every fixed
strip used above. Near terms use the distance bound; far differences
are $O((t-\gamma)^{-2})$ and are summed by the local count.
For $w=C/L\le1/10$, the Euler series at $1+w$ and the functional
equation at $-w$ also give the uniform convexity bound
\[
 |L_\chi(\sigma+it)|\ll L(qT)^{(1+w-\sigma)/2},
       \quad -w\le\sigma\le1+w,\quad T/2\le t\le3T.
\]
One direct maximum-principle proof applies it to
$L_\chi(s)(qT)^{(s-1-w)/2}e^{(s-it)^2}$.
The boundary estimates are absolute, and the Gaussian removes the
horizontal edges. For zeta, first multiply by $(s-1)/(s+1)$.
Slightly wider strips and Cauchy circles of radius $c/L$ give fixed
shift and derivative versions.

The exact main terms in \eqref{eq:dirichlet-three-residue},
\eqref{eq:dirichlet-discrete}, and \eqref{eq:bilinear-first-shift},
and the imaginary-part bound in \eqref{eq:real-bilinear}, consequently
have the following common uniform error:
\begin{equation}\label{eq:uniform-moment-error-new}
                         O_\epsilon(qT^{1/2+\epsilon}y).
\end{equation}
We verify every source of the conductor factor. Opposite $L$ factors
on a high horizontal edge have product $O(L^C\sqrt{qT})$ by the
preceding convexity bound. Their coefficient-polynomial product is
$O_\epsilon(y^{1+\epsilon}L^C)$, and the good-height logarithmic
derivative costs $O(L^2)$. A bounded bottom and its low occurrences
cost $O_\epsilon(qyT^\epsilon)$; the same strip proof applies with
height replaced by three. Endpoint restoration and $O(1/t)$ gamma
errors cost $O_\epsilon(\sqrt{qT}\,yT^\epsilon)$.

In the stationary transform the critical point is
$t=2\pi m/(qk)$ and the exact scalar is $\tau(\bar\chi)/q$.
The pre-transform prefactor is $q^{c-1/2}k^{c-1}m^{-c}$,
where $c=1+O(L^{-1})$. After the first and lower endpoint errors in
the cited stationary-phase lemma, it becomes $k^{-1/2}m^{-1/2}$.
Their sum is
\[
 k^{-1/2}\sum_{m\le CqkT}\frac{\tau_R(m)}{\sqrt m}
                    \ll_\epsilon\sqrt{qT}(qkT)^\epsilon.
\]
The upper endpoint error is bounded by the same prefactor times
$\sum_{m\asymp qkT}(|m-qkU/(2\pi)|+qk\sqrt T)^{-1}\ll L$.
This includes the nearest integers. The nonstationary ranges and
the relative Stirling error give $O(\sqrt{qT}L^C)$ per outer
coefficient by integration by parts and the absolute Dirichlet norm.
Restoring critical points below the bottom cutoff costs
$O_\epsilon(\sqrt q\,yT^\epsilon)$. Summing proves the required
stationary error, with no conductor in the amplitude constant.

Finally write $d_1=(n,k)$, $Q=qk/d_1$ in the auxiliary-character
average. The exact factorization and Parseval estimate already
proved give the uniform Perron integrand bound
\[
 (d_1Q(|\Im s|+3)T)^\epsilon\tau_R(d_1)\sqrt{Q+y}.
\]
Take Perron height $T^{10}$, a half-integer temporary cutoff, and
shift to $1/2+\eta$ with $\eta$ small compared with the final
exponent. The cutoff lies between $T/(2\pi)$ and $CT^{5/2}$.
Truncate after the character combination, whose original multiplier
has modulus one and whose nonprincipal remainder has modulus at
most two. Thus no character-count loss enters the truncation.
The remaining vertical error is at most
\[
 T^\epsilon(qTk/d_1)^{1/2}\tau_R(d_1)\sqrt{qk/d_1+y}.
\]
Its two outer sums are bounded by
\begin{align*}
 q\sqrt T\,T^\epsilon
 \sum_{k\le y}|a_2(k)|\sum_{d_1\mid k}\frac{\tau_R(d_1)}{d_1}
                       &\ll_\epsilon q\sqrt T\,yT^\epsilon,\\
 \sqrt{qT}\sqrt y\,T^\epsilon
 \sum_{k\le y}\frac{|a_2(k)|}{\sqrt k}
       \sum_{d_1\mid k}\frac{\tau_R(d_1)}{\sqrt{d_1}}
                       &\ll_\epsilon\sqrt{qT}\,yT^\epsilon.
\end{align*}
The beneficial outer factor $q^{-1/2}$ has even been discarded.
For the principal residue cutoff, let $\mathcal F(s)$ be the
principal Dirichlet series in the preceding Perron integral.
Enclose all its possible shifted poles by $|s-1|=C_2/L$, choosing
the fixed $C_2$ larger than all scaled shift bounds. On this circle
the finite polynomial has a fixed logarithmic Dirichlet norm,
and the valuation bounds and principal Euler factors give
\[
 |\mathcal F(s)|\ll_\epsilon
       \tau_R(d_1)(d_1Q)^\epsilon L^{C_3}.
\]
The residue sum and its derivative are
\[
 \mathcal R(x)=\frac1{2\pi i}\oint\mathcal F(s)\frac{x^s}{s}\,ds,
 \qquad
 \mathcal R'(x)=\frac1{2\pi i}\oint\mathcal F(s)x^{s-1}\,ds.
\]
Since $x\ll T^{5/2}$, the factor $x^{s-1}$ is bounded on this
circle. After decreasing preliminary exponents, this gives
$\mathcal R'(x)\ll_\epsilon T^\epsilon\tau_R(d_1)$.
Changing $x$ to $\lfloor x\rfloor+1/2$ therefore has this same
bound; the outer divisor sums put it within
\eqref{eq:uniform-moment-error-new}. Fixed shift derivatives add
only logarithmic powers.
Pole cutoff changes and horizontal Perron edges lie within these
bounds. The principal Gauss scalar and all residue identities remain
exactly real, so the signed cancellation has the same error.
Fixed Cauchy derivatives add logarithmic powers only. This proves
\eqref{eq:uniform-moment-error-new}.

\paragraph{The primitive Gauss factor and the error term.}
The same proof gives the sharper common error
\begin{equation}\label{eq:retained-gauss-moment-error}
                         O_\epsilon(\sqrt{qT}\,yT^\epsilon),
                 \qquad 1\le q\le T,
\end{equation}
in \eqref{eq:dirichlet-three-residue},
\eqref{eq:dirichlet-discrete}, \eqref{eq:bilinear-first-shift}
and \eqref{eq:real-bilinear}, with the same fixed divisor envelopes,
length bound $y\le T^\theta$, $\theta<1/2$, and fixed shift circles.
Indeed, $|\tau(\bar\chi)/q|=q^{-1/2}$. Retaining this factor in
the two outer sums just estimated gives
\begin{align*}
 \sqrt{qT}\,T^\epsilon
 \sum_{k\le y}|a_2(k)|\sum_{d_1\mid k}\frac{\tau_R(d_1)}{d_1}
                    &\ll_\epsilon\sqrt{qT}\,yT^\epsilon,\\
 \sqrt T\sqrt y\,T^\epsilon
 \sum_{k\le y}\frac{|a_2(k)|}{\sqrt k}
       \sum_{d_1\mid k}\frac{\tau_R(d_1)}{\sqrt{d_1}}
                    &\ll_\epsilon\sqrt T\,yT^\epsilon.
\end{align*}
The preliminary exponents have again been decreased before absorbing
the fixed divisor and logarithmic factors.

Here are the remaining cutoff checks for this stronger bound.
The Perron horizontal edges have integrand bounded by
$T^\epsilon\tau_R(d_1)\sqrt{Q+y}\,x/H$, whereas the vertical
bound uses $T^\epsilon\tau_R(d_1)\sqrt{Q+y}\sqrt x$.
Since $x\ll T^{5/2}$ and $H=T^{10}$, the horizontal edges lie within
the same two outer bounds. Combining characters before truncation
gives the half-integer Perron error
$O(T^\epsilon d_1^\epsilon x/H)$. Its outer sum, with the Gauss
factor retained, is
$O_\epsilon(\sqrt q\,T^{1+\epsilon}y/H)$.

For the principal residue restoration, the displayed contour for
$\mathcal R'(x)$ gives
$\mathcal R'(x)\ll_\epsilon T^\epsilon\tau_R(d_1)$.
The principal character coefficient has modulus at most one by
Parseval. Consequently changing the half-integer cutoff back to
$x$ costs at most
\[
 q^{-1/2}T^\epsilon
 \sum_{k\le y}\frac{|a_2(k)|}{k}
                  \sum_{d_1\mid k}\tau_R(d_1)
                    \ll_\epsilon q^{-1/2}T^\epsilon,
\]
which is also within \eqref{eq:retained-gauss-moment-error}.
The factors produced by shifts of size $O(L^{-1})$ are bounded here;
fixed shift derivatives contribute only logarithmic powers.

The high horizontal edges, stationary-phase errors, gamma remainder
and high endpoint restorations were already bounded by
$O_\epsilon(\sqrt{qT}\,yT^\epsilon)$ before the weaker common
bound was taken. The bounded bottom and its low occurrences cost
$O_\epsilon(qyT^\epsilon)$, and $q\le\sqrt{qT}$ on $1\le q\le T$.
The restored low stationary points cost
$O_\epsilon(\sqrt q\,yT^\epsilon)$; the right-edge harmonic
off-diagonal costs $O_\epsilon(yT^\epsilon)$. Both satisfy the
stronger bound. This proves \eqref{eq:retained-gauss-moment-error},
including fixed Cauchy derivatives. The later conductor ranges and
choice of $\kappa$ continue to use \eqref{eq:uniform-moment-error-new}.

\subsection{Deleted Euler factors and the exact rescaling}
Let $E_q=\prod_{p\mid q}(1-p^{-1})$. Splitting the prime divisors
at $\log(q+3)$, and using $\sum_{p\mid q}\log p\le\log q$, gives
\[
 \sum_{p\mid q}\frac1p=O(1+\log\log\log(q+10)),\qquad
 E_q^{-1}\ll1+\log\log(q+3).
\]
For $q\le T$, $|w|\le C/\sqrt L$, split the logarithmic derivative
of the deleted factors at $p=L^2$. The small-prime part is $O_C(V)$,
because $p^{|\Re w|}\le e^{2CV/\sqrt L}$; the large-prime part is
at most $CL^{-2+2C/\sqrt L}\sum_{p\mid q}\log p=O_C(1)$.
Integration from zero to $w$ proves
\begin{equation}\label{eq:uniform-euler-new}
                Z_q(1+w)=\frac{E_q}{w}(1+O_C(|w|V))
\end{equation}
uniformly, with the residue interpretation at zero.
The local remainder factors in the Euler products are
$1+O(p^{-3/2})$ in a common neighborhood and equal one at zero.
Omitting any subset of primes therefore preserves their uniform
boundedness and convergence to one by a common summable majorant.
On the full Mellin lines, the restricted Euler series bound both
$Z_q$ and its reciprocal by $\zeta(\Re s)$ when $\Re s>1$.
The global bounds $L^3,L^6,L^2$ consequently acquire at most a
fixed factor $V^B$. The order-sixteen Mellin weights still give
tails $O(V^BL^{-5})$, $O(V^BL^{-2})$, and $O(V^BL^{-6})$.
On the remaining box \eqref{eq:uniform-euler-new} has relative
error $O(V/\sqrt L)=o(1)$. All $E_q$ powers cancel before physical
shift differentiation, exactly as in the fixed-character formulas.

Put $\mathcal L=\log(qT)=\ell L$. The actual same polynomial has
length $T^\nu=e^{(\nu/\ell)\mathcal L}$ and cutoff denominator
$\nu L=(\nu/\ell)\mathcal L$.
Changing every Mellin and physical shift variable from units $L^{-1}$
to units $\mathcal L^{-1}$ therefore gives exactly the earlier
kernels with length parameter $\nu/\ell$.
The gamma logarithm is $\mathcal L+O(1)$ and
$N\sim T\mathcal L/(2\pi)$ uniformly. A physical derivative
introduces $\ell$, and two introduce $\ell^2$ in the normalizations
used here. Thus the two actual constants are
\begin{equation}\label{eq:uniform-moment-constants-new}
                  a_\ell=\ell a(P,\nu/\ell),\qquad
                  b_\ell=\ell^2b(P,\nu/\ell).
\end{equation}
This is the precise rescaling of the finite kernels.
Their integrable majorants also prove continuity in $\ell$ near one.
The certified strict gap at one permits an absolute $\kappa_0>0$
such that, for $1\le\ell\le1+\kappa_0$,
\[
 0<a_\ell<2,\quad0<b_\ell<2,\quad
 a_\ell^2-\frac{b_\ell}{2(1-1/2500)}>\frac1{4000}.
\]
Choose once $0<\kappa<\min(1/10,d/4,\kappa_0)$.
For every fixed insertion degree,
$qT^{1/2+\epsilon}T^{2\nu+o(1)}=o(T)$ uniformly when $\epsilon$
is sufficiently small. The undifferentiated moment and Cauchy proof
of \eqref{eq:weighted-curvature-new} now have absolute constants
and one uniform eventual height on $q\le T^\kappa$.

\subsection{Uniform Gaussian restrictions and continuation}
The five moment limits of Theorem~\ref{thm:gaussian-direction-moments}
are uniform in this conductor range, with $a,b$ replaced by
$a_\ell,b_\ell$. Here are the remaining error checks.
Deleting prime divisors of $q$ changes $S_T$ by $O(1+\log V)$, so
\[
 S_T=V-(\log V)^2-2\log\log V+O(1+\log V),\qquad
                      \sigma_T-s_T=o(1)
\]
uniformly. The rational Landau contour has uniform good-height,
gamma and Euler bounds for $q\le T$, so its error, all diagonal
and repeated-prime counts, and the $T^{o(1)}$ off-diagonal errors
are uniform. The marked local Euler coefficients remain the same
four polynomial identities. Their central-box error is a fixed
power of $V$ times $d_TRV+z_T^{-1}$, with $R=d_T^{-1/2}$.
The intermediate box costs $O(V^{B+K}R^{-14})$ and the central
Laurent relative error is $O(RV/L)$, both $o(1)$. Formula
\eqref{eq:uniform-moment-error-new} handles all finite polynomial
insertions and physical shift derivatives.

For the even Hardy phase, retaining $q$ in the proof of
\eqref{eq:rational-phase-new} gives the same exact main coefficient
and the uniform error
\[
 O_\epsilon\left(T^{1/2+\epsilon}\sqrt{qb}
          +(q^{-1}a/b+q^{-1/2}\sqrt{a/b})L^2\right).
\]
The $q/\varphi(q)$ in its main term is $O(V)$.
After the exterior truncation and fixed prime insertions, the
main term is $O_k(T(1+V)^{k+1})$ and the errors are bounded by
$\sqrt{qT}T^{\epsilon+o(1)}$, $q^{-1}T^{1/4+o(1)}$ and
$q^{-1/2}T^{o(1)}$. They are $o(N)$ uniformly after fixing a small
$\epsilon$. For the signed moment choose
$\delta_0>0$ with $\kappa+\nu+\delta_0<1/2$ and use
\eqref{eq:uniform-moment-error-new} at length
$T^{\nu+\delta_0+o(1)}$; the exterior tail is removed by the
already uniform second and prime moments. Thus all five polynomial
limits are uniform. The interval conclusion follows from the same
Taylor-remainder and Fourier argument. Explicitly, a failure of
uniformity would give a sequence with a convergent subsequence of
$\ell\in[1,1+\kappa]$; the uniform moments and continuity in
\eqref{eq:uniform-moment-constants-new} give the required interval
limit on that subsequence, a contradiction.

The logarithm estimate \eqref{eq:log-kernel-new} is uniform: its
constant has no conductor, trivial zeros depend only on parity, and
its contour now has the uniform good-height bounds.
The coefficient comparison has squared norm
$O(1+(\log V)^2+\log V)=o(V)$ uniformly.
The reversed kernel mean is also uniform, since the ambient count
in $[T/2,3T]$ is $O(T\log(qT))=O(N)$ and the outer shell counts
are $O(2^jT(L+j))$. Thus every $o(N)$ derivative-selection
exception and its $o(NL)$ signed mass is uniform.

We finish with explicit margins. On $I=[1,2]$, of mass $p>1/24$,
the uniform moment gap gives
$\sum_{I,|\cos\psi|\ge1/50}|A_\rho|\ge(p/20000-o(1))NL$.
Hence total absolute $C_\rho$ mass is at least
$(p/3000000-o(1))NL$, and each sign has at least
$pNL/10000000$ eventually. We have
$\sum|A_\rho|^2\le3NL^2$ and $|C_\rho|\le6|A_\rho|$.
Take $\omega=10^{-11}$ and phase margin $10^{-14}$.
Endpoint strips cost $o(NL)$, the phase deletion at most
$6\sqrt3\,10^{-14}NL$, and choose $K_0$ so large that
$6\sqrt{C_M}/K_0<\omega/10$ for the curvature deletion.
The uniform logarithm argument gives the derivative floor
$L e^{\sigma_T/2-1}/(2\eta)$ on all but $o(N)$ of these separated
centers. Each sign still has mass $\omega NL$; its squared mass
is at most $108NL^2$. Thus it has at least
$\omega^2N/108>10^{-30}N$ centers.
All local deformation constants are absolute because their exterior
Dirichlet series have coefficients bounded by one. The proofs of
Theorems~\ref{thm:common-deformation} and
\ref{thm:anisotropic-continuation} therefore give the first
assertion with a uniform height and constants. The logarithm of the
left side on $|u|\le\sqrt V/(100T)$ is at most
$0.12\sqrt V+O(\log V)$, below $s_T+\log c_1$ eventually.

For an arbitrary fixed interval of mass $p>0$, repeat the restricted
selection with its initial sign mass $pNL/10000000$.
Choose its curvature cutoff so the loss is below $pNL/10^9$.
The uniform transfer of logarithmic derivatives leaves sign mass
$pNL/10^8$ with squared mass at most $144pNL^2$.
Cauchy gives at least $pN/(144\cdot10^{16})>10^{-30}pN$ centers.
Inner intervals and the exact reflected speed identities prove the
uniform distribution assertions. Applying this selection to
$[A+1,A+2]$ gives $|D_\rho|\ge Le^{As_T}$, after $o(N)$
deletions, and ample slack above the stated $10^{-31}$ fraction.
The same absolute difference estimate
\eqref{eq:anisotropic-perturbation} and the same relative direction
bound prove the final region and real-ray conclusions. This
completes every asserted uniformity statement.
\end{proof}
